% !PRINTED 48
\weberpage{48}

\begin{center}
\S.\,12.

\gesperrt{Divisoren einer Abel'schen Gruppe}.

\gesperrt{Reciproke Gruppen}.
\end{center}

Ein Theiler $T$ einer \gesperrt{Abel}'schen Gruppe $S$ ist, wie schon oben bemerkt, immer ein Normaltheiler, und daher giebt es nach \S.\,4 eine zu $T$ complementäre Gruppe
\[
S/T,
\]
deren Elemente die Nebengruppen von $T$ sind, nämlich
\weq{1}{T,\ TA',\ TA''\ldots,}
worin $A', A'' \ldots$ gewisse Elemente aus $S$ bedeuten. Die Anzahl der Elemente (1), also der Grad der Gruppe $S/T$ ist, wenn $t$ der Grad von $T$ ist, gleich dem Index des Theilers $T$ von $S$, also gleich
\[
\frac{n}{t} = j.
\]

Diese Gruppe $S/T$ ist aber selbst wieder eine \gesperrt{Abel}'sche; denn es ist
\[
TA'\,TA'' = TA'A'',\ \ TA''\,TA' = TA''A',
\]
also beides einander gleich.

Es sind nun die Charaktere der Gruppe $S/T$ zu bestimmen. Wir bezeichnen die Elemente dieser Gruppe mit
\weq{2}{T,\ T_1,\ T_2 \ldots T_{j-1},}
und einen ihrer Charaktere mit $\xi(T_i)$.

Aus dieser Function $\xi(T)$ können wir nun eine Function $\xi(A)$ der Elemente von $S$ ableiten, indem wir
\weq{3}{\xi(A_i) = \xi(T_i)}
setzen, wenn $A_i$ irgend ein in der Nebengruppe $T_i$ vorkommendes Element ist. Für die Elemente der Gruppe $T$ selbst ist dann $\xi(A) = 1$.

Diese Function $\xi(A_i)$ ist aber unter den Charakteren von $S$ enthalten. Denn wenn $A_i$ in $T_i$ und $A_k$ in $T_k$ vorkommt, so kommt $A_iA_k$ in $T_iT_k$ vor, und folglich ist
\weq{4}{\xi(A_i)\xi(A_k) = \xi(T_i)\xi(T_k) = \xi(T_iT_k) = \xi(A_iA_k).}
Dies aber ist nach \S.\,11,(4) die Definition für einen Charakter von $S$.
